\scanpage{48}

On va définir aussi
\[
\Delta_I\in A(I\amalg I)
\]
qui sera tel que
\[
\widetilde{\unsure{\delta|\Delta}}_I(x)=x\qquad(x\in A(I))
\qquad\text{i.e.}\quad\widetilde{\unsure{\delta|\Delta}}_I=\operatorname{id}_{A(I)}.
\]
Pour le voir, on \deleted{définit} considère $\delta:I\amalg I-I$, d'où
\[
\delta^{\deleted{\unsure{*}}}_*:A(I)-A(I\amalg I),
\]
on pose
\[
\Delta_I=\delta_*(1)
\]
On aura (\unsure{écrivant} $\Delta$ au lieu de $\Delta_I$
\margin{
\diagram{
\begin{tikzcd}
& X^I \arrow[d,no head,"\delta"] &\\
& X^{I\amalg I} \arrow[dl,"p_1"'] \arrow[dr,"p_2"] &\\
X^I && X^I
\end{tikzcd}
}
}
\[
\begin{aligned}
\widetilde\Delta(x)&\overset{\text{déf}}{=}p_{2*}\bigl(\Delta\,p_1^*(x)\bigr)\\
&=p_{2*}\bigl(\underbrace{\delta_*(1)\,p_1^*(x)}
_{\substack{\Vert\\ \unsure{\text{pr.}}\ \delta_*(\delta^*p_1^*(x))=\delta_*(x)}}\bigr)\\
&=p_{2*}\delta_*(x)=x
\end{aligned}
\]

Il faut encore \unsure{définir|vérifier} l'associativité de l'opération $\bullet$, et le fait que les $\unsure{\delta|\Delta}_I$ jouent le rôle d'unités. On voit que sous ces conditions,
\margin{47}
